\pdfbookmark[0]{Anzeigen. Untersuchungen über Themen der höheren Geodäsie. 1843.}{anzeigen-geodaesie-erste}
\begin{center}
\rule{.85\textwidth}{.3pt}\par
Göttingische gelehrte Anzeigen.\quad 6. November 1843.\par
\rule{.85\textwidth}{.3pt}
\end{center}
\hspace*{\parindent}Am 23. Oktober überreichte Hofrat \textsc{Gauss} der Königlichen Sozietät einen Vortrag mit dem Titel
\begin{center}\textit{Untersuchungen über Themen der höheren Geodäsie},\end{center}
\noindent über den hier kurz berichtet werden soll.

Beim trigonometrischen Teil der Gradmessung, die der Verfasser in den Jahren 1821---1827 ausführte, und bei den späteren damit zusammenhängenden trigonometrischen Vermessungen, die sich über das ganze Königreich Hannover erstreckten, wurden andere Wege eingeschlagen als sonst üblich. Das gilt für die Art, wie die Messungen vorgenommen wurden, und noch mehr für ihre anschließende mathematische Behandlung und die Gewinnung der Ergebnisse. Einige der für Hofrat \textsc{Gauss} charakteristischen Methoden sind zwar bereits bekannt gemacht worden, teils von ihm selbst in verschiedenen schon vor Langem erschienenen Aufsätzen, teils von anderen, die nach mündlichen oder brieflichen Mitteilungen bei ihren eigenen trigonometrischen Messungen davon Gebrauch gemacht hatten. Doch der bedeutendere Teil dieser Methoden ist bisher noch nicht dargestellt worden. Es sind diejenigen, die sich am stärksten von den sonst gebräuchlichen unterscheiden und für deren Verständnis eine tiefere mathematische Begründung erforderlich ist. Der Verfasser hatte früher vor, nach dem vollständigen Abschluss der Messungen diese selbst zusammen mit allen von ihm angewandten Verfahren in einem eigenen Werk darzulegen. Umstände, deren Erläuterung nicht hierhergehört, haben dieses Vorhaben vorerst aufgeschoben. Deshalb hat er sich dafür entschieden, in einer Reihe einzelner Abhandlungen zu veröffentlichen, was an der %
